\chapter{争议怎样改变周秦史的读法}

本附录不是把每一处疑问裁决为一个标准答案，也不重述全卷的争议账本。以下十例之所以被选出，是因为若把后出的故事、精确数字或现代概念当作同时代事实，读者对王朝建立、霸政、变法、统一和覆亡的理解都会随之改变。每一例先分开材料所处的层次，再给出本卷可以依凭的说法、不能越过的界线，以及尚待材料缩小的空白。

\section{西周：从胜利、摄政到王室危机}

\subsection{牧野：一个甲子日，不是一份战报}

\eventref{WZ-1046-MUYE}{西周·牧野}是周取代商的重大断裂，但其绝对年代和战场情景不能混为同一问题。利簋是接近事件的金文物证；《尚书·牧誓》是传世政治文本；《史记·周本纪》是汉代回望；王年、日食等综合推算则属于现代年代解释。它们彼此照亮，却各自回答不同的问题。

\paragraph{可以这样读。}前1046年是现行的综合估算；利簋所见“甲子朝、征商”使这个时间锚点不可忽略。金文与传世传统共同支撑周代商兴亡这一历史断裂。

\paragraph{不可越过。}不能把前1046年写成全无争议的精确年份，也不能由甲子日复原战场地点、兵数、阵势、人物言辞或一份完整战报；后出的正当化叙事不能直接充作同时代记录。

\paragraph{尚待材料。}出土背景明确、可以互校的早周铭文，能重复检验的年代模型，以及与商周转折地层相连的材料，或可收窄年代；它们仍未必补得上战场细节。

\subsection{周公摄政：相对次序，不等于逐年程序}

武王死后、东方危机、\eventref{WZ-1040-DONGZHENG}{西周·三监东征}和归政的传统，常被读成一份可以逐年排开的政务记录。《金縢》《大诰》《康诰》等《尚书》篇章有复杂的成文、传本和政治语境；《史记》把它们整合进王朝叙事；大盂鼎等金文能旁证处分、册命和政治网络，却没有留下完整年表。

\paragraph{可以这样读。}幼主继承、东方危机、严厉处置及其后东方关系的重组，构成可靠的相对次序；周公摄政与归政也是早期周政治传统的重要部分。

\paragraph{不可越过。}摄政的精确年数、每篇《尚书》是否为当日诰命原件、管叔、蔡叔、武庚的动机，以及每次战事的部署，都不能当作已证事实。

\paragraph{尚待材料。}可断代且出土背景可靠的周初王命、军旅铭文和早期写本的文本比较，能够缩小用语和顺序的范围；一件器物仍不能替代整段政治史。

\subsection{共和：可靠的纪年锚点，未定的执政形式}

\eventref{WZ-0841-GONGHE}{西周·共和}在传统连续纪年中提供前841年的稳定锚点，也指向厉王出奔后的非常继承危机。可是《国语·周语上》及《史记·周本纪》的周、召共和叙事，与《竹书纪年》所存共伯和“干王位”的异说，不能被压成一种已经由同代材料裁定的制度。

\paragraph{可以这样读。}共和对应王权危机中的非通常继承安排，前841年可用作连续纪年的锚点。

\paragraph{不可越过。}“共和”不是现代共和政体的名称；周、召与共伯和两种说法都不能被写成唯一确定的执政者或体制，也不能据此把此前所有西周年份装作同等精确。

\paragraph{尚待材料。}可核王年的晚西周铭文、不同《竹书纪年》系统的文本史研究和王系交叉考订，或可解释两种叙事的关系；直接行政文书的缺失仍是根本限制。

\subsection{镐京失守：不能只读成宫廷报应}

\eventref{WZ-0771-HAOJING}{西周·镐京陷落}常被后出兴亡故事收束为幽王个人失德、褒姒与烽火的因果链。《国语》《左传》《史记》保存的叙事并不相同；《诗经·小雅》和晚西周金文能提供危机时代的环境和关系线索，但不是陷落当天的记录。

\paragraph{可以这样读。}幽王死、镐京陷落和东迁构成可靠的年代骨架。继承冲突、申侯联盟、边防压力和关中权力网络的重组，是解释危机时不可略去的结构条件。

\paragraph{不可越过。}褒姒不笑、反复举烽、诸侯受骗的连锁故事，没有早期材料的独立证实；“犬戎”等名称也不能被压缩为一个固定的现代族群或唯一攻城主体。

\paragraph{尚待材料。}关中晚西周聚落、破坏层、铭文的可比断代，以及传世文本的早期异文，可帮助辨认危机的时序与参与者；宫廷传奇很难由这类材料验证。

\section{春秋与战国：会盟和分期不能替代政治过程}

\subsection{召陵与葵丘：齐的动员能力，不是对楚的属国统治}

\eventref{SA-0656-SHAOLING}{春秋·召陵之盟}与\eventref{SA-0651-KUIQIU}{春秋·葵丘之会}确认齐能组织跨国行动，却不等于齐获得了常设帝国。 《春秋》经文保留会、伐、盟等最小行动；《左传》展开屈完、苞茅和昭王不返的交涉；《国语·齐语》与《史记》再把它们放入齐霸兴亡线。没有已知盟书或行政簿册可将谈判补成逐字记录。

\paragraph{可以这样读。}前656年齐方联盟到楚北并以盟收束，前651年齐侯盟诸侯于葵丘；两事显示齐有跨国调动能力，也显示楚仍是不能忽略的独立政治体。

\paragraph{不可越过。}楚不能因此被写成齐的属国；使者问答、盟辞全貌和周室授命礼仪不可当作原始记录；一次会盟也不等于齐能够常设征税、任命或裁决诸侯。

\paragraph{尚待材料。}可确证的盟誓实物、同代铭文，以及路线和驻军的区域考古，有助于检验参与者和行动范围；地理材料本身不能证明谈判措辞或意图。

\subsection{春秋何时结束：文本终点、权力重组和名分追认}

\eventref{SA-0481-LIN}{春秋·获麟与陈氏政变}所在的前481年、\eventref{WQ-0453-ZHI}{战国·三家灭智}所在的前453年和\eventref{WQ-0403-THREEJIN}{战国·三家分晋}所在的前403年，划出的并不是同一种边界。《春秋》停止于前481年是经文事实；《史记》世家、六国年表与《资治通鉴》组织三家分晋和册命；把它们称作“春秋”或“战国”的界线，则是现代解释的选择。

\paragraph{可以这样读。}前481年是传世鲁国经文的终点；前453年显示晋国卿族权力的实质重组；前403年是周王室对韩、赵、魏既成地位的名分承认。

\paragraph{不可越过。}没有任何一年可以被写成全国制度一夜突变的唯一开端；册命不能倒推三家此前没有国家能力，《春秋》停笔也不代表列国政治同时停止。

\paragraph{尚待材料。}可交叉断代的封君、官署、兵器和城邑材料，可让区域变化的速度更清楚；它们不能消除“以何种变化为界”的解释选择。

\section{战国：改革与战争的尺度}

\subsection{商鞅：改革的节点，不是一人写完的制度}

\eventref{WQ-0356-SHANGYANG}{战国·商鞅变法}、迁咸阳与第二次改革是秦国制度路径上的可靠节点，但改革清单、文本归属和成熟制度之间有距离。《史记·商君列传》是汉代传记；《商君书》有复杂的成书和篇章层次；睡虎地、里耶简是较晚且地点有限的法律、县廷文书；现代制度史才将它们与《秦本纪》交读。

\paragraph{可以这样读。}秦孝公任用商鞅、改革与迁咸阳的编年主线可靠。军功、户籍、基层控制、县邑与尺度的重组形成了秦可持续发展的制度路径，商鞅死后新法并未随之废止。

\paragraph{不可越过。}《商君书》的每一句都不能当作前356年的诏令；后来的成熟秦律、县制和度量衡也不能全部归功于一人。改革不是无阻力、一次完成的“现代化”。

\paragraph{尚待材料。}早于帝国期、出处可靠的秦法律和行政简牍，以及可断代的官印、权量器、城址，可检验制度的初行与扩展；局部文书仍不能证明全国同步。

\subsection{长平：决定性失败，不能化约为一个数字或一个人物}

\eventref{WQ-0260-CHANGPING}{战国·长平之战}使赵的野战能力和战略位置遭受严重损失，必须从上党争夺、长期相持、秦的补给围困与赵的政治决策一起理解。《史记·白起王翦列传》《赵世家》及《资治通鉴》保留基本战事，也保存高度戏剧化的操作细节；高平一带战国遗存只是考古旁证，不能换算出史书总数。

\paragraph{可以这样读。}前260年赵军遭遇决定性失败；秦的补给与围困、赵的决策和资源约束都是解释结局的必要条件。

\paragraph{不可越过。}“四十余万”不是已核军籍，每名降卒的死亡方式和埋葬时序也不能断言；赵括不能被故事化为唯一的“纸上谈兵”原因，长平更不能被直接等同于前228年赵亡。

\paragraph{尚待材料。}经合规发掘、分层发表的人骨和兵器材料，结合道路、仓储和人口承载研究，可讨论规模区间和战场过程；没有独立军籍，精确伤亡数仍不可得。

\section{秦：统一的过程与覆亡的多重条件}

\subsection{统一六国：有编年骨架，不等于预设总计划}

\eventref{QIN-0230-UNIFICATION}{秦·统一六国}可以确认韩、赵、魏、楚、燕、齐在前230至前221年间先后失去主体政权。可是《史记·秦始皇本纪》《六国年表》和各国世家提供的是后出的编年骨架，《资治通鉴》再作串联；并没有六场连续的同代军政档案。战场地理、水利和交通研究只能检验行动条件。

\paragraph{可以这样读。}秦逐步利用前次接收的道路、粮食和人力，同时面对各国不同的地理与联盟困局。都城失守、残余政权和地方接管须分别理解：赵都失守或燕失蓟，都不等于领土立即尽平。

\paragraph{不可越过。}王翦请兵“六十万”、大梁水攻的工期、燕辽东路线、齐廷劝说和每城降附人数，不能写成无争议事实；六国覆亡的先后也不能证明秦以一套不变总计划和同一种战术横扫天下。

\paragraph{尚待材料。}出土背景可靠的郡县文书、封泥、兵器和城址破坏或续存层，连同可复核的水利和交通调查，可细化接管过程；它们不能直接证明宫廷预先设计过一套总计划。

\subsection{秦亡：继承失信、动员压力与指挥链断裂}

从\eventref{QIN-0210-SUCCESSION}{秦·沙丘继承}到\eventref{QIN-0209-CHENSHENG}{秦·大泽起事}，再到秦主力崩解与前206年子婴降服，秦亡有一条牢靠的时间主线。《史记》的《陈涉世家》《项羽本纪》《高祖本纪》《秦始皇本纪》是汉代回望中的人物与兴亡叙事；《资治通鉴》提供后出的逐年框架；睡虎地、里耶简可说明期限、征发、连带责任与县廷操作，却不能验证大泽乡的每一情节。

\paragraph{可以这样读。}前210年的继承危机、前209年的反秦动员、前207年秦主力崩解和前206年的降服相互衔接；继承失信、战争消耗、地方倒戈与指挥链断裂是叠加的条件。

\paragraph{不可越过。}赵高、暴雨、严刑或巨鹿之战，都不能被写成唯一亡秦原因。迟到必死、鱼腹书、破釜沉舟、酷刑对白和赵高一人亡秦，也不应当作无媒介的事实；秦律的局部条文不能直接变成大泽乡的逐项原因。

\paragraph{尚待材料。}有年号的秦末地方文书、交通与军需材料，以及经严格发表的战场遗存，能改善对地方响应和军队构成的判断；失传的中央军政档案仍无法由它们替代。

\begin{turningpoint}[读争议，不是把历史读成空白]
这十例所留下的并非“什么都不知道”。它们保住了周代商、周公传统、共和危机、镐京陷落、齐的会盟、晋的权力重组、秦的改革与战争、统一及其覆亡的年代骨架和政治后果；同时也要求读者把金文、经年、传世叙事、后设史书、简牍考古与现代解释放回各自能够承担的范围。最有力量的叙事，不是把空白填成细节，而是让材料已经支持的行动、约束与代价清楚地显现出来。
\end{turningpoint}
